% The main argument uses integral observations of the two-mode model.
% General statistical bounds and detailed interval estimates remain in supplement.tex.
\subsection{From state observations to analytic error}
\label{core-sec-learning-kam}
Long-time prediction tests the fitted dynamics along selected trajectories.
To establish invariant motion, we also need error bounds on a neighborhood
of those trajectories. Following the learning-to-KAM connection of
\citet[Theorem~11]{chen2022}, we derive these bounds for the analytic
readout~\eqref{eq:unified-feature-readout},
$h_\eta=h_2+\eta^{\mathsf T}F$, whose quadratic part is fixed and whose
full-cosine target is $h_\star=h_{\eta_\star}$. We write $\widehat h$ for
the fitted function, including coefficient-export rounding.

For independently sampled gradient labels, our uniform learning bound
in Supplementary Theorem~\ref{S-thm:learning-to-cl} connects nonzero loss
to a $C^\ell$ error under explicit sampling and regularity assumptions.
Supplementary Corollary~\ref{S-cor:common-target-kam} combines it with
\citet[Theorem~A]{bounemoura}. The COMSOL observations used here are
correlated states along trajectories, so we instead derive a deterministic
bound from their integrated equations of motion.

COMSOL supplies $17$ position--velocity observations on each of $m=64$
conservative two-mode trajectories. Let $c_j:[0,T]\to\mathbb R^4$ be
their degree-eight fitted curves, with endpoints fixed to the observations,
and let $Lx=J\nabla h_2(x)$. Define
\begin{equation}
 \begin{aligned}
  D_j&=T^{-1}\int_0^T J\nabla F(c_j(t))\,dt,\\
  y_j&=T^{-1}\left[c_j(T)-c_j(0)-\int_0^T Lc_j(t)\,dt\right],\\
  G_I&=(4m)^{-1}\sum_jD_j^{\mathsf T}D_j.
 \end{aligned}
 \label{core-eq-integral-design}
\end{equation}
These labels require states and feature integrals, without acceleration
reconstruction. Let $r_I$ be the RMS residual $D_j\widehat\eta-y_j$
over the $4m$ scalar observations, and let $\epsilon_I$ bound the RMS
integrated physical defect
$e_j=T^{-1}\int_0^T(\dot c_j-J\nabla h_\star(c_j))\,dt$.

\begin{proposition}[Integral observations bound the learned energy]
\label{core-prop-integral-bound}
If $\lambda_{\min}(G_I)\geq\underline\lambda>0$, then
\begin{equation}
 \begin{aligned}
  \|\widehat\eta-\eta_\star\|_2
   &\leq b_I:=\frac{r_I+\epsilon_I}{\sqrt{\underline\lambda}},\\
  \|\widehat h-h_\star\|_r
   &\leq b_I B_F(r)+\delta_{h,\mathrm{dec}}(r).
 \end{aligned}
 \label{core-eq-integral-bound}
\end{equation}
Here $\|\cdot\|_r$ is the supremum norm on the complex polydisc
$|x|_\infty<r$, $B_F(r)^2=\sum_j\|F_j\|_r^2$, and
$\delta_{h,\mathrm{dec}}$ bounds coefficient-export rounding.
\end{proposition}
\begin{proof}
Endpoint matching gives $y_j-D_j\eta_\star=e_j$.
The triangle inequality in the normalized stacked observation norm
bounds $\sqrt{\underline\lambda}\|\widehat\eta-\eta_\star\|_2$
by $r_I+\epsilon_I$. Cauchy--Schwarz in feature space then gives
the analytic bound.
\end{proof}
The physical defect is evaluated against the declared target; it is
additional information beyond the training loss. Interval integration
bounds this defect and retains the full sine remainder and coefficient
rounding. Cauchy's estimate then bounds every derivative on a smaller
polydisc. Supplementary Section~\ref{S-sec:validated-integral-learning}
gives these estimates and the independent trajectory-error bound.

\subsection{The magnetic nonlinearity preserves twist after learning}
\label{core-sec-twist}
The analytic error bound becomes useful for KAM theory because the
magnetic potential supplies a nondegenerate frequency map.
For $h_\star=h_2+g(d^{\mathsf T}X)$, with
$g(u)=1-\cos u-u^2/2$, its quartic part is
$-(d^{\mathsf T}X)^4/24$. Averaging over the modal angles with
$I_j=(X_j^2+Y_j^2)/2$ gives the action Hessian
\begin{equation}
 \begin{aligned}
 A_\star&=-\frac18
 \begin{pmatrix}
 d_1^4&2d_1^2d_2^2\\2d_1^2d_2^2&d_2^4
 \end{pmatrix},\\
 \det A_\star&=-\frac3{64}d_1^4d_2^4\ne0.
 \end{aligned}
 \label{core-eq-magnetic-twist}
\end{equation}
Indeed, $\langle X_j^4\rangle=3I_j^2/2$ and
$\langle X_1^2X_2^2\rangle=I_1I_2$, while odd terms average to zero.
Both modal tip rotations are nonzero, so the twist is nonsingular,
although indefinite.

Learning preserves this property if its quartic error is smaller than
the target's least singular value. Let $B$ map the readout coefficients
to quartic coefficients and let $\delta_{\mathrm{dec}}$ bound rounding
in that map. Angular averaging yields
\begin{equation}
 \|A_{\widehat h}-A_\star\|_2
 \leq\sqrt{19}\bigl(\|B\|_2b_I+\delta_{\mathrm{dec}}\bigr)
 <\sigma_{\min}(A_\star).
 \label{core-eq-twist-transfer}
\end{equation}
The factor $\sqrt{19}$ bounds the quartic-to-twist operator, including
momentum-dependent terms. Table~\ref{core-tab-transfer} establishes
the strict inequality for the Hamiltonian learned from the COMSOL states.

\begin{table}[t]
\centering\small
\caption{Learning error and twist margin for the COMSOL two-mode fit.
All displayed bounds include coefficient rounding and are rounded outward.}
\label{core-tab-transfer}
\begin{tabular}{lr}
\toprule
Quantity & Bound\\
\midrule
Training MSE & $[3.6030,3.6031]\,10^{-23}$\\
$r_I$, $\epsilon_I$ & $\leq6.003\,10^{-12},\ 6.884\,10^{-12}$\\
$\lambda_{\min}(G_I)$ & $\geq2.930\,10^{-7}$\\
$b_I$ & $\leq2.381\,10^{-8}$\\
Twist error & $\leq1.038\,10^{-7}$\\
$\sigma_{\min}(A_\star)$ & $\geq4.903\,10^{-4}$\\
$\|\widehat h-h_\star\|_{0.3}$ & $\leq1.331\,10^{-9}$\\
\bottomrule
\end{tabular}
\end{table}

The complete magnetic potential also contributes a sixth-order analytic
tail bounded by $(Dr)^6\cosh(Dr)/720$, where $D=\|d\|_1$.
Together with fourth-order nonresonance, the twist bound gives a local
normal form $\nu\cdot I+\tfrac12I^{\mathsf T}A_{\widehat h}I+O(r^6)$.
Our remainder estimates in Supplementary
Sections~\ref{S-sec:unified-normal-form}--\ref{S-sec:unified-learning-twist}
therefore yield positive-measure KAM families near equilibrium.
To establish persistence at a specified physical amplitude, we next
apply an a posteriori theorem to the same fitted function.

\subsection{An invariant torus at finite physical amplitude}
\label{core-sec-finite-torus}
The existence theorem used here is
\citet[Theorem~3.1]{haro2019aposteriori}, with the constants of
arXiv:1711.05515v1, Appendix~B. Our contribution in the following result
is the construction and analytic bounds for the learned magnetic model.

\begin{theorem}[Finite-amplitude persistence for the learned magnetic model]
\label{core-thm-learned-torus}
The fitted Hamiltonian in Table~\ref{core-tab-transfer}, specified by its
exact coefficients in Supplementary
Section~\ref{S-sec:validated-integral-learning}, admits a real analytic
embedded Lagrangian torus $K_\infty:\mathbb T^2\to\mathbb R^4$ satisfying
\begin{equation}
 J\nabla\widehat h(K_\infty)=DK_\infty\varpi.
 \label{core-eq-torus-frequency}
\end{equation}
The exact frequency $\varpi$ is specified in
Supplementary Eqs.~\eqref{S-eq:certified-invariance}--\eqref{S-eq:certified-quadratic-frequency};
angles have period one and time is dimensionless.
The torus is within $2.0894\,10^{-41}$ of the computed Fourier embedding
on the complex strip of width $1/80$. Its beam reconstruction attains
a tip displacement of at least $1.8977904$ mm and has centreline
displacement at most $1.8977962$ mm at every phase.
\end{theorem}
\begin{proof}
The bounded continued-fraction coefficients of the frequency ratio give
$|k\cdot\varpi|\geq(9/1750)|k|_1^{-1}$ for every nonzero integer $k$,
and hence the weaker Diophantine exponent $3$ used in the theorem.
On the strip $\rho=1/40$, 512-bit ball-arithmetic evaluation of the Fourier candidate,
including interpolation tails, bounds the analytic invariance residual
by $2.7310\,10^{-77}$ and its inverse averaged torsion by $80526$.
The geometric and derivative estimates in Supplementary
Section~\ref{S-sec:unified-aposteriori} give a residual threshold greater
than $8.6086\,10^{-76}$, with smallness quotient below $0.031724<1$.
The cited theorem then supplies the exact torus and correction bound.
Fourier bounds and the Hermite-element convex-hull property give the
displacement estimates.
\end{proof}
Thus nonzero learning loss, preserved twist and finite-amplitude
persistence are established for one fitted Hamiltonian. The result uses
the independently simulated two-mode state observations; the full-beam
comparisons in Section~\ref{sec:analysis} quantify a separate reduction
error. Section~\ref{sec:application} connects this conservative motion
to the controlled full beam.

\subsection{From one torus to a family of quasiperiodic motions}
\label{core-sec-positive-measure}
The nondegeneracy bound survives the a posteriori iteration:
$\|\langle T_\infty\rangle^{-1}\|\leq2\,10^5$ for the exact torus's
averaged torsion. In an analytic symplectic tube, solving one quadratic
homological equation therefore gives
\begin{equation}
 \widehat h\circ\Psi
 =e+\varpi\cdot I+\tfrac12I^{\mathsf T}\overline A I+O(|I|^3),
 \qquad\det\overline A\ne0.
 \label{core-eq-torus-normal-form}
\end{equation}
Our construction identifies $\overline A$ with the averaged torsion and
retains an analytic remainder (Supplementary
Corollary~\ref{S-cor:finite-torus-families}). The nonresonant frequency,
nondegenerate quadratic part and analytic remainder meet the hypotheses
of \citet[Theorem~6, $\alpha=1$, arXiv:1412.0509v1]{bounemoura2016torus}.
That theorem gives positive-measure families of analytic Lagrangian
invariant tori around $K_\infty$, with relative density tending to one
as the transverse action box shrinks. The same argument applies to all
seven fitted or target Hamiltonians considered in the supplement.

A numerical fraction additionally depends on the chosen initial-state
domain. For $U\subset\mathbb R^4$ of positive Liouville volume, define
$f_{\rm KAM}(U)=\mu_L(U\cap\mathcal K_h)/\mu_L(U)$, where
$\mathcal K_h$ is the union of analytic Diophantine Lagrangian tori.
Supplementary Section~\ref{S-app:measure-proof} applies
\citet[Theorems~1 and~3]{chierchia2021measure} to an explicit tube;
its very small transverse width limits its use as an operating tolerance.
For the broader action domain $I_1,I_2\in[0.0032,0.0072]$, independent
uniform sampling of actions and phases gives $1024/1024$ finite-time
quasiperiodicity candidates over $8192$ dimensionless time units
($717.5$ s). At step $1/64$, all pass the eight-digit test on consecutive windows
of length $2048$ and $4096$; six- and ten-digit tests agree.
Supplementary Section~\ref{S-sec:measure-numerical} gives the weighted
Birkhoff method, uncertainty and refinement checks. This broader-domain
diagnostic complements the local family theorem without supplying an
infinite-time volume bound. The magnetic realization below carries the
family's low-mode parameter measure into the damped plant; this measure
is distinct from volume in the full $128$-dimensional state space.
